\documentclass[UTF8,fontset=fandol,11pt]{ctexart}
\usepackage[a4paper,margin=24mm]{geometry}
\usepackage{amsmath,amssymb,amsthm,graphicx,tabularx,booktabs,hyperref}
\newtheorem{definition}{定义}[section]
\newtheorem{theorem}[definition]{定理}
\newtheorem{example}[definition]{例题}
\newtheorem{exercise}[definition]{练习}
\hypersetup{unicode=true,colorlinks=true,linkcolor=blue}
\title{最小二乘法讲义计划}
\date{}
\begin{document}
\maketitle
\section{一元最小二乘拟合：从平方误差到正规方程}\label{sec:S01}
% UNIT S01.U01
\begin{definition}
设有 $n\geq 2$ 个观测点 $(x_i,y_i)$，其中 $x_i,y_i\in\mathbb{R}$，$i=1,\ldots,n$。用直线
\[
\hat y_i=a+bx_i
\]
拟合被解释变量 $y_i$，其中 $a,b\in\mathbb{R}$ 分别为截距参数和斜率参数。定义平方误差目标函数为
\begin{equation}
S(a,b)=\sum_{i=1}^{n}\bigl[y_i-(a+bx_i)\bigr]^2.
\label{def:ols-objective}
\end{equation}
一元最小二乘估计 $(\hat a,\hat b)$ 是使 $S(a,b)$ 在 $\mathbb{R}^2$ 上达到最小值的参数对。
\end{definition}
% UNIT S01.U02
\begin{definition}
对给定的参数 $a,b$，第 $i$ 个观测的残差定义为观测值与拟合值之差：
\begin{equation}
e_i=y_i-(a+bx_i).
\label{def:residual}
\end{equation}
当参数取最小二乘估计 $(\hat a,\hat b)$ 时，记
\[
\hat y_i=\hat a+\hat b x_i,
\qquad
\hat e_i=y_i-\hat y_i.
\]
此时平方误差目标函数的最小值为残差平方和 $\sum_{i=1}^n\hat e_i^2$。
\end{definition}
% UNIT S01.U03
\begin{theorem}
设 $x_i,y_i\in\mathbb{R}$，且
\label{thm:quadratic-unique-minimum}
\[
\sum_{i=1}^{n}(x_i-\bar x)^2>0,
\qquad
\bar x=\frac{1}{n}\sum_{i=1}^{n}x_i,
\qquad
\bar y=\frac{1}{n}\sum_{i=1}^{n}y_i.
\]
则平方误差函数 $S(a,b)$ 在 $\mathbb{R}^2$ 上具有唯一最小点 $(\hat a,\hat b)$。该最小点满足正规方程
\begin{equation}
\sum_{i=1}^{n}\bigl[y_i-(\hat a+\hat b x_i)\bigr]=0,
\qquad
\sum_{i=1}^{n}x_i\bigl[y_i-(\hat a+\hat b x_i)\bigr]=0,
\label{thm:normal-equations}
\end{equation}
并且
\[
\hat b=\frac{\sum_{i=1}^{n}(x_i-\bar x)(y_i-\bar y)}{\sum_{i=1}^{n}(x_i-\bar x)^2},
\qquad
\hat a=\bar y-\hat b\bar x.
\]
\end{theorem}
% UNIT S01.U04
\begin{proof}
将 $S(a,b)$ 展开得
\[
S(a,b)=\sum_{i=1}^{n}\bigl[y_i-a-bx_i\bigr]^2.
\]
对 $a$ 和 $b$ 分别求偏导，得到
\[
\frac{\partial S}{\partial a}=-2\sum_{i=1}^{n}\bigl[y_i-a-bx_i\bigr],
\qquad
\frac{\partial S}{\partial b}=-2\sum_{i=1}^{n}x_i\bigl[y_i-a-bx_i\bigr].
\]
因此，最小点必须满足式 \eqref{thm:normal-equations}。第一条正规方程给出
\[
n\hat a+\hat b\sum_{i=1}^{n}x_i=\sum_{i=1}^{n}y_i,
\]
从而
\[
\hat a=\bar y-\hat b\bar x.
\]
把它代入第二条正规方程，并利用
\[
\sum_{i=1}^{n}(x_i-\bar x)=0,
\qquad
\sum_{i=1}^{n}(y_i-\bar y)=0,
\]
可得
\[
\hat b\sum_{i=1}^{n}(x_i-\bar x)^2
=\sum_{i=1}^{n}(x_i-\bar x)(y_i-\bar y).
\]
由于分母严格为正，故 $\hat b$ 唯一确定，随后 $\hat a$ 也唯一确定。

下面验证由上述公式给出的 $(\hat a,\hat b)$ 确实满足正规方程。对第一条，有
\[
\sum_{i=1}^{n}\bigl[y_i-(\hat a+\hat b x_i)\bigr]
=\sum_{i=1}^{n}y_i-n\hat a-\hat b\sum_{i=1}^{n}x_i
=n\bar y-n(\bar y-\hat b\bar x)-\hat b n\bar x=0.
\]
对第二条，利用 $\sum_{i=1}^{n}(x_i-\bar x)=0$ 和 $\sum_{i=1}^{n}(y_i-\bar y)=0$，有
\[
\sum_{i=1}^{n}x_i\bigl[y_i-(\hat a+\hat b x_i)\bigr]
=\sum_{i=1}^{n}(x_i-\bar x)(y_i-\bar y)
-\hat b\sum_{i=1}^{n}(x_i-\bar x)^2=0,
\]
最后一步由 $\hat b$ 的定义得到。因此 $(\hat a,\hat b)$ 满足式 \eqref{thm:normal-equations}。

为说明这是唯一的全局最小点，令 $u=a-\hat a$，$v=b-\hat b$。由于 $(\hat a,\hat b)$ 满足正规方程，直接展开平方项可得
\[
S(a,b)=S(\hat a,\hat b)+\sum_{i=1}^{n}(u+vx_i)^2.
\]
再将 $x_i=(x_i-\bar x)+\bar x$ 代入，并令 $w=u+v\bar x$，得到
\[
\sum_{i=1}^{n}(u+vx_i)^2
=nw^2+v^2\sum_{i=1}^{n}(x_i-\bar x)^2\geq 0.
\]
当且仅当 $w=0$ 且 $v=0$ 时等号成立；由此 $u=0$。所以 $(\hat a,\hat b)$ 是唯一最小点。
\end{proof}
% UNIT S01.U05
\begin{theorem}
在定理 \eqref{thm:normal-equations} 的条件下，最小二乘拟合直线经过样本均值点 $(\bar x,\bar y)$，即
\[
\hat y(\bar x)=\hat a+\hat b\bar x=\bar y.
\]
等价地，含截距的一元最小二乘拟合满足
\[
\sum_{i=1}^{n}\hat e_i=0.
\]
\end{theorem}
% UNIT S01.U06
\begin{proof}
由定理 \eqref{thm:normal-equations} 给出的截距公式，
\[
\hat a=\bar y-\hat b\bar x.
\]
因此
\[
\hat a+\hat b\bar x=\bar y.
\]
另一方面，第一条正规方程为
\[
\sum_{i=1}^{n}\hat e_i=\sum_{i=1}^{n}\bigl[y_i-(\hat a+\hat b x_i)\bigr]=0\text{，}
\]
故拟合残差的总和为零。
\end{proof}
% UNIT S01.U07
\begin{example}
某地区三年的广告支出 $x$ 与销售额 $y$ 的观测数据为
\[
(x_1,y_1)=(1,2),\qquad (x_2,y_2)=(2,3),\qquad (x_3,y_3)=(3,5)\text{。}
\]
这三个观测点的样本均值为
\[
\bar x=\frac{1+2+3}{3}=2,\qquad \bar y=\frac{2+3+5}{3}=\frac{10}{3}\text{。}
\]
用一元直线拟合销售额与广告支出之间的线性关系。
\end{example}
% UNIT S01.U08
\par\noindent\textbf{解：}样本均值为
\[
\bar x=2,
\qquad
\bar y=\frac{10}{3}.
\]
因此
\[
\sum_{i=1}^{3}(x_i-\bar x)(y_i-\bar y)
=(-1)\left(-\frac43\right)+0\left(-\frac13\right)+1\left(\frac53\right)=3,
\]
且
\[
\sum_{i=1}^{3}(x_i-\bar x)^2=(-1)^2+0^2+1^2=2\text{。}
\]
所以
\[
\hat b=\frac{3}{2},
\qquad
\hat a=\frac{10}{3}-\frac32\cdot2=\frac13\text{。}
\]
拟合直线为
\[
\hat y=\frac13+\frac32x\text{。}
\]
三个观测点的拟合值和残差分别为
\[
\begin{array}{c|c|c}
 i&\hat y_i&\hat e_i=y_i-\hat y_i\\ \hline
 1&\frac{11}{6}&\frac16\\
 2&\frac{10}{3}&-\frac13\\
 3&\frac{29}{6}&\frac16
\end{array}
\]
可见 $\hat e_1+\hat e_2+\hat e_3=0$，且拟合直线经过 $(\bar x,\bar y)=(2,10/3)$。在经济解释上，该直线描述样本内广告支出与销售额的线性拟合关系，但不能仅凭拟合结果断言广告支出对销售额具有因果影响。\par
% UNIT S01.U09
\begin{exercise}
\label{mathagent:exercise:41706b9e371cbf8283a4f1192e660b4426644ffb8f650af52f47cef5491b1ca0}
设三家企业的培训支出 $x$ 与劳动生产率 $y$ 的数据为
\[
(x_1,y_1)=(1,4),\qquad (x_2,y_2)=(2,5),\qquad (x_3,y_3)=(4,8)\text{。}
\]
完成下列任务：
\begin{enumerate}
\item 计算 $\bar x$ 和 $\bar y$；
\item 求最小二乘拟合直线 $\hat y=\hat a+\hat b x$；
\item 计算三个观测点的残差，并验证残差之和为零；
\item 求培训支出为 $x=3$ 时的拟合值，并说明这是样本区间内的插值。
\end{enumerate}
\end{exercise}
\section{多元最小二乘与矩阵表示}\label{sec:S02}
% UNIT S02.U01
\begin{definition}
\label{def:design-matrix}设有 $n\geq 1$ 个观测，每个观测包含一个被解释变量 $y_i\in\mathbb{R}$ 和 $p-1$ 个解释变量 $x_{i1},\ldots,x_{i,p-1}\in\mathbb{R}$。设计矩阵 $X\in\mathbb{R}^{n\times p}$ 定义如下：当 $p\geq 2$ 时，
\[
X=\begin{bmatrix}
1&x_{11}&\cdots&x_{1,p-1}\\
1&x_{21}&\cdots&x_{2,p-1}\\
\vdots&\vdots&\ddots&\vdots\\
1&x_{n1}&\cdots&x_{n,p-1}
\end{bmatrix};
\]
当 $p=1$ 时，模型仅含截距项，$X=\begin{bmatrix}1&\cdots&1\end{bmatrix}^{\mathsf T}\in\mathbb{R}^{n\times 1}$。在后一情形中，约定空和 $\sum_{j=1}^{0}(\cdot)=0$。第一列的 $1$ 表示模型含有截距项。令参数列向量、响应列向量分别为
\[
\beta=\begin{bmatrix}\beta_0\\\beta_1\\\vdots\\\beta_{p-1}\end{bmatrix}\in\mathbb{R}^{p},
\qquad
y=\begin{bmatrix}y_1\\y_2\\\vdots\\y_n\end{bmatrix}\in\mathbb{R}^{n}.
\]
矩阵形式的线性拟合值为 $X\beta$，第 $i$ 个拟合值为 $\beta_0+\sum_{j=1}^{p-1}\beta_jx_{ij}$；特别地，当 $p=1$ 时，该式按上述空和约定退化为 $\beta_0$。
\end{definition}
% UNIT S02.U02
\begin{definition}
对任意参数向量 $\beta\in\mathbb{R}^{p}$，残差向量定义为
\begin{equation}
e(\beta)=y-X\beta.
\label{def:residual-vector}
\end{equation}
其第 $i$ 个分量为
\[
e_i(\beta)=y_i-\left(\beta_0+\sum_{j=1}^{p-1}\beta_jx_{ij}\right).
\]
矩阵形式的平方误差目标函数定义为
\[
Q(\beta)=e(\beta)^{\mathsf T}e(\beta)=(y-X\beta)^{\mathsf T}(y-X\beta)=\sum_{i=1}^{n}e_i(\beta)^2.
\]
当 $p=2$ 且 $X$ 的第二列为 $(x_1,\ldots,x_n)^{\mathsf T}$ 时，$Q(\beta)$ 就是前文的一元目标函数 $S(a,b)$，其中 $\beta=(a,b)^{\mathsf T}$。
\end{definition}
% UNIT S02.U03
\begin{definition}
若设计矩阵 $X\in\mathbb{R}^{n\times p}$ 的列向量线性无关，则称 $X$ 列满秩，即 $\operatorname{rank}(X)=p$。这要求 $p\leq n$，并且不存在非零向量 $v\in\mathbb{R}^{p}$ 使得 $Xv=0$。在列满秩条件下，矩阵 $X^{\mathsf T}X\in\mathbb{R}^{p\times p}$ 可逆。
\end{definition}
% UNIT S02.U04
\begin{theorem}
\label{thm:unique-least-squares-solution}设 $X\in\mathbb{R}^{n\times p}$，$y\in\mathbb{R}^{n}$，且 $X$ 列满秩，即 $\operatorname{rank}(X)=p$。定义 $Q(\beta)=(y-X\beta)^{\mathsf T}(y-X\beta)$。则 $Q$ 在 $\mathbb{R}^{p}$ 上存在唯一最小点 $\hat\beta$，并且该最小点满足矩阵正规方程
\begin{equation}
X^{\mathsf T}X\hat\beta=X^{\mathsf T}y.
\label{thm:matrix-normal-equation}
\end{equation}
由于 $X^{\mathsf T}X$ 可逆，唯一最小二乘解为
\[
\hat\beta=(X^{\mathsf T}X)^{-1}X^{\mathsf T}y.
\]
对应的拟合值和残差分别为 $\hat y=X\hat\beta$ 与 $\hat e=y-X\hat\beta$。
\end{theorem}
% UNIT S02.U05
\begin{proof}
将目标函数展开为
\[
Q(\beta)=y^{\mathsf T}y-2\beta^{\mathsf T}X^{\mathsf T}y+\beta^{\mathsf T}X^{\mathsf T}X\beta.
\]
对 $\beta$ 的各个分量求偏导，得到梯度
\[
\nabla Q(\beta)=-2X^{\mathsf T}y+2X^{\mathsf T}X\beta.
\]
因此，最小点必须满足 $\nabla Q(\hat\beta)=0$，也就是
\[
X^{\mathsf T}X\hat\beta=X^{\mathsf T}y,
\]
这就是式 \eqref{thm:matrix-normal-equation}。

由于 $X$ 列满秩，对任意非零向量 $v\in\mathbb{R}^{p}$，都有 $Xv\neq 0$，从而
\[
v^{\mathsf T}X^{\mathsf T}Xv=(Xv)^{\mathsf T}(Xv)=\sum_{i=1}^{n}(Xv)_i^2>0.
\]
所以 $X^{\mathsf T}X$ 可逆，正规方程至多有一个解，并且该解为
\[
\hat\beta=(X^{\mathsf T}X)^{-1}X^{\mathsf T}y.
\]

还需说明该解确实给出全局最小值。令 $h=\beta-\hat\beta$，则 $\beta=\hat\beta+h$，并且
\[
y-X\beta=(y-X\hat\beta)-Xh.
\]
展开平方误差可得
\[
Q(\beta)=Q(\hat\beta)-2h^{\mathsf T}X^{\mathsf T}(y-X\hat\beta)+h^{\mathsf T}X^{\mathsf T}Xh.
\]
由正规方程，$X^{\mathsf T}(y-X\hat\beta)=0$，因此
\[
Q(\beta)=Q(\hat\beta)+h^{\mathsf T}X^{\mathsf T}Xh
=Q(\hat\beta)+(Xh)^{\mathsf T}(Xh)\geq Q(\hat\beta).
\]
当且仅当 $Xh=0$ 时等号成立。列满秩又保证 $Xh=0$ 只能推出 $h=0$，即 $\beta=\hat\beta$。所以 $\hat\beta$ 是唯一的全局最小点。
\end{proof}
% UNIT S02.U06
\begin{theorem}
在定理 \eqref{thm:matrix-normal-equation} 的条件下，最小二乘残差 $\hat e=y-X\hat\beta$ 满足
\[
X^{\mathsf T}\hat e=0.
\]
也就是说，残差向量与设计矩阵的每一列都正交。特别地，当模型含截距项时，$X$ 的第一列为全 $1$ 向量，于是
\[
\sum_{i=1}^{n}\hat e_i=0.
\]
当模型还包含解释变量列 $(x_{1j},\ldots,x_{nj})^{\mathsf T}$ 时，还满足
\[
\sum_{i=1}^{n}x_{ij}\hat e_i=0,
\qquad j=1,\ldots,p-1.
\]
\end{theorem}
% UNIT S02.U07
\begin{proof}
由定理 \eqref{thm:matrix-normal-equation} 的正规方程，
\[
X^{\mathsf T}X\hat\beta=X^{\mathsf T}y.
\]
移项得到
\[
X^{\mathsf T}(y-X\hat\beta)=0.
\]
根据残差向量的定义 \eqref{def:residual-vector}，这就是
\[
X^{\mathsf T}\hat e=0.
\]
矩阵乘积的第 $j$ 个分量等于设计矩阵第 $j$ 列与残差向量的内积，因此每一列都与残差向量正交。若第一列为全 $1$ 向量，则对应内积为 $\sum_{i=1}^{n}\hat e_i$；若第 $j+1$ 列为解释变量 $x_j$，则对应内积为 $\sum_{i=1}^{n}x_{ij}\hat e_i$。
\end{proof}
% UNIT S02.U08
\begin{example}
某企业有三次观测，$y$ 表示劳动生产率，$x_1$ 表示培训支出，$x_2$ 表示设备投入。观测数据为
\[
\begin{array}{c|c|c|c}
 i&x_{i1}&x_{i2}&y_i\\ \hline
 1&1&0&4\\
 2&0&1&3\\
 3&1&1&5
\end{array}
\]
考虑含截距的模型
\[
y_i=\beta_0+\beta_1x_{i1}+\beta_2x_{i2}+e_i.
\]
写出设计矩阵 $X$、响应向量 $y$，并利用矩阵正规方程求最小二乘估计。
\end{example}
% UNIT S02.U09
\par\noindent\textbf{解：}由数据可得
\[
X=\begin{bmatrix}1&1&0\\1&0&1\\1&1&1\end{bmatrix},
\qquad
 y=\begin{bmatrix}4\\3\\5\end{bmatrix}.
\]
计算得到
\[
X^{\mathsf T}X=
\begin{bmatrix}3&2&2\\2&2&1\\2&1&2\end{bmatrix},
\qquad
X^{\mathsf T}y=
\begin{bmatrix}12\\9\\8\end{bmatrix}.
\]
因此正规方程为
\[
\begin{bmatrix}3&2&2\\2&2&1\\2&1&2\end{bmatrix}
\begin{bmatrix}\hat\beta_0\\\hat\beta_1\\\hat\beta_2\end{bmatrix}
=
\begin{bmatrix}12\\9\\8\end{bmatrix}.
\]
直接检验可得
\[
\hat\beta=
\begin{bmatrix}2\\2\\1\end{bmatrix},
\]
因为
\[
\begin{bmatrix}3&2&2\\2&2&1\\2&1&2\end{bmatrix}
\begin{bmatrix}2\\2\\1\end{bmatrix}
=
\begin{bmatrix}12\\9\\8\end{bmatrix}.
\]
所以拟合方程为
\[
\hat y=2+2x_1+x_2.
\]
三个观测的拟合值依次为 $4,3,5$，因此残差向量为
\[
\hat e=y-X\hat\beta=
\begin{bmatrix}0\\0\\0\end{bmatrix}.
\]
在这个例子中，三个观测点恰好被模型完全拟合。若给定一个新的解释变量取值 $(x_1,x_2)=(2,1)$，则将其代入拟合方程可得样本外预测值
\[
\widehat y(2,1)=2+2\times 2+1=7.
\]
由于样本中 $x_1$ 只取过 $0$ 和 $1$，这里的预测涉及超出样本观测范围的外推，应谨慎解释其可靠性。这个计算说明矩阵正规方程可以同时处理多个解释变量，但拟合关系仍然只是样本内的线性拟合，不能据此断言培训支出或设备投入对劳动生产率具有因果影响。\par
% UNIT S02.U10
\begin{exercise}
\label{mathagent:exercise:49ae28350c6a4abd6936597716c61d2e3cb07c4268424bce9420505d640e7b34}
设三次观测的设计矩阵和响应向量为
\[
X=\begin{bmatrix}1&0&1\\1&1&0\\1&1&1\end{bmatrix},
\qquad
 y=\begin{bmatrix}2\\4\\5\end{bmatrix}.
\]
完成下列任务：
\begin{enumerate}
\item 计算 $X^{\mathsf T}X$ 和 $X^{\mathsf T}y$；
\item 利用正规方程求 $\hat\beta=(\hat\beta_0,\hat\beta_1,\hat\beta_2)^{\mathsf T}$；
\item 写出拟合方程 $\hat y=\hat\beta_0+\hat\beta_1x_1+\hat\beta_2x_2$；
\item 计算残差向量，并验证它与 $X$ 的每一列正交。\end{enumerate}
\end{exercise}
\section{残差性质、拟合优度与经济解释}\label{sec:S03}
% UNIT S03.U01
\begin{definition}
在沿用上一节列满秩设计矩阵 $\operatorname{rank}(X)=p$ 的设定下，考虑含截距的最小二乘拟合。记第 $i$ 个观测的拟合值为 $\hat y_i$，残差为 $\hat e_i=y_i-\hat y_i$。若设计矩阵的第一列为全 $1$ 向量，则由正规方程有
\[
\sum_{i=1}^{n}\hat e_i=0.
\]
进一步，若某个向量与残差向量的内积为零，则称残差向量与该向量正交。特别地，在上述含截距且列满秩的模型中，残差向量与全 $1$ 向量正交；若模型还含有解释变量列，则残差向量也与每一个解释变量列正交。\label{def:orthogonal-residual}
\end{definition}
% UNIT S03.U02
\begin{definition}
设观测值为 $y_1,\ldots,y_n\in\mathbb{R}$，其样本均值为
\[
\bar y=\frac{1}{n}\sum_{i=1}^{n}y_i.
\]
总平方和定义为观测值相对于样本均值的离差平方和：
\[
\operatorname{TSS}=\sum_{i=1}^{n}(y_i-\bar y)^2.
\]
残差平方和定义为
\[
\operatorname{RSS}=\sum_{i=1}^{n}\hat e_i^2.
\]
总平方和衡量观测值围绕样本均值的总变异，残差平方和衡量拟合后尚未解释的平方变异。\label{def:total-sum-of-squares}
\end{definition}
% UNIT S03.U03
\begin{definition}
当模型含截距且 $\operatorname{TSS}>0$ 时，决定系数定义为
\[
R^2=1-\frac{\operatorname{RSS}}{\operatorname{TSS}}.
\]
由此，$R^2$ 表示相对于以样本均值作为拟合值的基准，平方误差减少的比例。若 $\operatorname{RSS}\leq\operatorname{TSS}$，则 $0\leq R^2\leq 1$；但 $R^2$ 高并不表示解释变量与被解释变量之间存在因果关系。\label{def:coefficient-of-determination}
\end{definition}
% UNIT S03.U04
\begin{theorem}
设 $n\geq 1$，$y_i,\hat y_i\in\mathbb{R}$，并令
\[
\hat e_i=y_i-\hat y_i,
\qquad
\bar y=\frac{1}{n}\sum_{i=1}^{n}y_i.
\]
若拟合模型含截距，并且残差满足
\[
\sum_{i=1}^{n}\hat e_i=0,
\qquad
\sum_{i=1}^{n}\hat e_i(\hat y_i-\bar y)=0,
\]
则有平方和分解
\begin{equation}
\sum_{i=1}^{n}(y_i-\bar y)^2
=
\sum_{i=1}^{n}(\hat y_i-\bar y)^2
+
\sum_{i=1}^{n}\hat e_i^2.
\label{thm:sum-of-squares-decomposition}
\end{equation}
因此，当 $\operatorname{TSS}>0$ 时，决定系数也可写为
\[
R^2=
\frac{\sum_{i=1}^{n}(\hat y_i-\bar y)^2}
{\sum_{i=1}^{n}(y_i-\bar y)^2}.
\]
\end{theorem}
% UNIT S03.U05
\begin{proof}
由 $y_i=\hat y_i+\hat e_i$，有
\[
y_i-\bar y=(\hat y_i-\bar y)+\hat e_i.
\]
对等式两边平方并求和，得到
\[
\begin{aligned}
\sum_{i=1}^{n}(y_i-\bar y)^2
={}&\sum_{i=1}^{n}(\hat y_i-\bar y)^2
+2\sum_{i=1}^{n}\hat e_i(\hat y_i-\bar y)
+\sum_{i=1}^{n}\hat e_i^2.
\end{aligned}
\]
根据假设
\[
\sum_{i=1}^{n}\hat e_i(\hat y_i-\bar y)=0,
\]
交叉项消失，于是得到式 \eqref{thm:sum-of-squares-decomposition}。

下面沿用上一节的列满秩条件 $\operatorname{rank}(X)=p$，并假定模型含截距。在这一条件下，残差与拟合值的内积也为零。事实上，\(\hat y\) 是设计矩阵列向量的线性组合，而由正规方程得到的残差与设计矩阵的每一列正交，因此
\[
\sum_{i=1}^{n}\hat e_i\hat y_i=0.
\]
再结合 \(\sum_{i=1}^{n}\hat e_i=0\)，可得
\[
\sum_{i=1}^{n}\hat e_i(\hat y_i-\bar y)
=\sum_{i=1}^{n}\hat e_i\hat y_i-\bar y\sum_{i=1}^{n}\hat e_i=0.
\]
所以平方和分解适用于上述列满秩的含截距最小二乘拟合。最后，将分解式代入
\[
R^2=1-\frac{\operatorname{RSS}}{\operatorname{TSS}}
\]
即可得到决定系数的第二种表达式。
\end{proof}
% UNIT S03.U06
\begin{example}
考虑三年广告支出与销售额数据
\[
(x_1,y_1)=(1,2),\qquad (x_2,y_2)=(2,3),\qquad (x_3,y_3)=(3,5).
\]
由一元最小二乘拟合可得
\[
\hat y=\frac{1}{3}+\frac{3}{2}x,
\]
相应的拟合值为
\[
\hat y_1=\frac{11}{6},\qquad
\hat y_2=\frac{10}{3},\qquad
\hat y_3=\frac{29}{6}.
\]
计算总平方和、残差平方和、解释平方和以及决定系数，并说明决定系数的经济含义及其解释边界。
\end{example}
% UNIT S03.U07
\par\noindent\textbf{解：}样本均值为
\[
\bar y=\frac{2+3+5}{3}=\frac{10}{3}.
\]
因此总平方和为
\[
\begin{aligned}
\operatorname{TSS}
&=\left(2-\frac{10}{3}\right)^2
+\left(3-\frac{10}{3}\right)^2
+\left(5-\frac{10}{3}\right)^2\\
&=\left(-\frac{4}{3}\right)^2
+\left(-\frac{1}{3}\right)^2
+\left(\frac{5}{3}\right)^2
=\frac{14}{3}.
\end{aligned}
\]
残差为
\[
\hat e_1=2-\frac{11}{6}=\frac{1}{6},
\qquad
\hat e_2=3-\frac{10}{3}=-\frac{1}{3},
\qquad
\hat e_3=5-\frac{29}{6}=\frac{1}{6}.
\]
它们的和为
\[
\hat e_1+\hat e_2+\hat e_3
=\frac{1}{6}-\frac{1}{3}+\frac{1}{6}=0.
\]
残差平方和为
\[
\operatorname{RSS}
=\left(\frac{1}{6}\right)^2
+\left(-\frac{1}{3}\right)^2
+\left(\frac{1}{6}\right)^2
=\frac{1}{6}.
\]
由平方和分解，解释平方和为
\[
\operatorname{ESS}=\operatorname{TSS}-\operatorname{RSS}
=\frac{14}{3}-\frac{1}{6}=\frac{9}{2}.
\]
所以决定系数为
\[
R^2=1-\frac{1/6}{14/3}
=1-\frac{1}{28}
=\frac{27}{28}.
\]
这表示在该样本中，广告支出与销售额的线性拟合相对于只用样本均值预测销售额，减少了 $27/28$ 的平方误差。这个数值描述的是样本内的拟合程度，不能单独证明广告支出对销售额具有因果影响，也不能保证对样本区间之外的数据同样准确。

还应注意异常观测的影响：若某一年销售额明显偏离总体趋势，它可能对斜率、拟合直线和残差平方和产生较大影响，从而显著改变 $R^2$。因此，较高的 $R^2$ 不一定表示所有观测都被同样良好地解释，分析时应检查是否存在离群观测，并结合数据背景判断其原因。\par
% UNIT S03.U08
\begin{exercise}
\label{mathagent:exercise:80c9d88bd59912a783f0955435a3f123c4b9f6c751a875a01f8d2dcd6b734e5e}
某企业三年的培训支出与劳动生产率观测值为
\[
(x_1,y_1)=(1,1),\qquad (x_2,y_2)=(2,2),\qquad (x_3,y_3)=(3,4).
\]
一位同学声称含截距的一元最小二乘拟合直线为
\[
\hat y=-\frac{1}{6}+\frac{3}{2}x.
\]
请检查这一说法，并用数据重新计算正确的最小二乘直线。然后完成下列任务：
\begin{enumerate}
\item 计算 $\bar y$、三个拟合值和三个残差；
\item 计算 $\operatorname{TSS}$、$\operatorname{RSS}$ 和 $R^2$；
\item 利用平方和分解验证 $\operatorname{TSS}=\operatorname{ESS}+\operatorname{RSS}$；
\item 说明该拟合结果在经济解释上的一个限制。
\end{enumerate}
\end{exercise}
\section*{练习答案}
% UNIT S01.U10
\par\noindent\textbf{\hyperref[mathagent:exercise:41706b9e371cbf8283a4f1192e660b4426644ffb8f650af52f47cef5491b1ca0]{练习 \ref*{mathagent:exercise:41706b9e371cbf8283a4f1192e660b4426644ffb8f650af52f47cef5491b1ca0} 的解答：}}答案如下：
\[
\bar x=\frac{1+2+4}{3}=\frac73,
\qquad
\bar y=\frac{4+5+8}{3}=\frac{17}{3}\text{。}
\]
进一步有
\[
\sum_{i=1}^{3}(x_i-\bar x)(y_i-\bar y)=\frac{19}{3},
\qquad
\sum_{i=1}^{3}(x_i-\bar x)^2=\frac{14}{3}\text{。}
\]
所以
\[
\hat b=\frac{19/3}{14/3}=\frac{19}{14},
\qquad
\hat a=\frac{17}{3}-\frac{19}{14}\cdot\frac73=\frac52\text{。}
\]
拟合直线为
\[
\hat y=\frac52+\frac{19}{14}x\text{。}
\]
三个拟合值和残差为
\[
\begin{array}{c|c|c}
 i&\hat y_i&\hat e_i=y_i-\hat y_i\\ \hline
 1&\frac{27}{7}&\frac17\\
 2&\frac{73}{14}&-\frac{3}{14}\\
 3&\frac{111}{14}&\frac{1}{14}
\end{array}
\]
残差之和为
\[
\frac17-\frac{3}{14}+\frac{1}{14}=0\text{。}
\]
最后，$x=3$ 时的拟合值为
\[
\hat y(3)=\frac52+\frac{19}{14}\cdot3=\frac{46}{7}\text{。}
\]
由于 $3$ 位于样本解释变量的最小值 $1$ 与最大值 $4$ 之间，这是样本区间内的插值。\par
% UNIT S02.U11
\par\noindent\textbf{\hyperref[mathagent:exercise:49ae28350c6a4abd6936597716c61d2e3cb07c4268424bce9420505d640e7b34]{练习 \ref*{mathagent:exercise:49ae28350c6a4abd6936597716c61d2e3cb07c4268424bce9420505d640e7b34} 的解答：}}先计算矩阵乘积：
\[
X^{\mathsf T}X=
\begin{bmatrix}3&2&2\\2&2&1\\2&1&2\end{bmatrix},
\qquad
X^{\mathsf T}y=
\begin{bmatrix}11\\9\\7\end{bmatrix}.
\]
正规方程为
\[
\begin{bmatrix}3&2&2\\2&2&1\\2&1&2\end{bmatrix}
\begin{bmatrix}\hat\beta_0\\\hat\beta_1\\\hat\beta_2\end{bmatrix}
=
\begin{bmatrix}11\\9\\7\end{bmatrix}.
\]
解得
\[
\hat\beta=
\begin{bmatrix}1\\3\\1\end{bmatrix}.
\]
因此拟合方程为
\[
\hat y=1+3x_1+x_2.
\]
拟合值为
\[
X\hat\beta=
\begin{bmatrix}2\\4\\5\end{bmatrix}=y,
\]
故残差向量为
\[
\hat e=y-X\hat\beta=
\begin{bmatrix}0\\0\\0\end{bmatrix}.
\]
于是
\[
X^{\mathsf T}\hat e=
\begin{bmatrix}0\\0\\0\end{bmatrix},
\]
这验证了残差向量与设计矩阵的每一列正交。\par
% UNIT S03.U09
\par\noindent\textbf{\hyperref[mathagent:exercise:80c9d88bd59912a783f0955435a3f123c4b9f6c751a875a01f8d2dcd6b734e5e]{练习 \ref*{mathagent:exercise:80c9d88bd59912a783f0955435a3f123c4b9f6c751a875a01f8d2dcd6b734e5e} 的解答：}}样本均值为
\[
\bar y=\frac{1+2+4}{3}=\frac{7}{3}.
\]
先计算
\[
\bar x=\frac{1+2+3}{3}=2,
\qquad
S_{xx}=(-1)^2+0^2+1^2=2,
\qquad
S_{xy}=(-1)\left(-\frac{4}{3}\right)+0\left(-\frac{1}{3}\right)+1\left(\frac{5}{3}\right)=3.
\]
因此正确的最小二乘斜率和截距分别为
\[
\hat b=\frac{S_{xy}}{S_{xx}}=\frac{3}{2},
\qquad
\hat a=\bar y-\hat b\bar x
=\frac{7}{3}-\frac{3}{2}\cdot2=-\frac{2}{3}.
\]
所以正确的拟合直线为
\[
\hat y=-\frac{2}{3}+\frac{3}{2}x.
\]
题目中给出的截距 $-1/6$ 不正确；其对应残差和不为零，因而不可能是这组含截距数据的最小二乘直线。

由正确的拟合直线，三个拟合值为
\[
\hat y_1=-\frac{2}{3}+\frac{3}{2}=\frac{5}{6},
\qquad
\hat y_2=-\frac{2}{3}+3=\frac{7}{3},
\qquad
\hat y_3=-\frac{2}{3}+\frac{9}{2}=\frac{23}{6}.
\]
相应残差为
\[
\hat e_1=1-\frac{5}{6}=\frac{1}{6},
\qquad
\hat e_2=2-\frac{7}{3}=-\frac{1}{3},
\qquad
\hat e_3=4-\frac{23}{6}=\frac{1}{6},
\]
且
\[
\hat e_1+\hat e_2+\hat e_3=0.
\]
总平方和为
\[
\operatorname{TSS}
=\left(1-\frac{7}{3}\right)^2
+\left(2-\frac{7}{3}\right)^2
+\left(4-\frac{7}{3}\right)^2
=\frac{14}{3}.
\]
残差平方和为
\[
\operatorname{RSS}
=\left(\frac{1}{6}\right)^2
+\left(-\frac{1}{3}\right)^2
+\left(\frac{1}{6}\right)^2
=\frac{1}{6}.
\]
解释平方和为
\[
\operatorname{ESS}=\operatorname{TSS}-\operatorname{RSS}
=\frac{14}{3}-\frac{1}{6}=\frac{9}{2}.
\]
因此确有
\[
\operatorname{TSS}=\operatorname{ESS}+\operatorname{RSS}
=\frac{9}{2}+\frac{1}{6}=\frac{14}{3}.
\]
决定系数为
\[
R^2=1-\frac{\operatorname{RSS}}{\operatorname{TSS}}
=1-\frac{1/6}{14/3}
=\frac{27}{28}.
\]
这说明该样本内的线性拟合相对于用样本均值预测，减少了 $27/28$ 的平方误差；但它不能单独证明培训支出对劳动生产率具有因果作用，也不能保证样本区间之外的预测准确。若存在异常观测，斜率、残差平方和和 $R^2$ 也可能受到较大影响。\par
\end{document}
